\hypertarget{data}{}
\hypertarget{data-portfolio-construction-units-and-provenance}{%
\chapter{Data, portfolio construction, units and
provenance}\label{data-portfolio-construction-units-and-provenance}}

This chapter fixes the raw material of the whole study: the analytical
grain of a single observation, the units in which returns and factors are
measured, how the 25 test portfolios and the factors are constructed, the
fail-closed checks that a panel must pass before any regression is run,
and why reconciling a file against its official source is a different
question from reproducing a number from that file. The estimator of
Chapter~\ref{preliminaries} is only as trustworthy as the panel it
consumes, so the discipline here is not clerical --- it is part of the
statistical specification.

\hypertarget{the-5-5-test-asset-grid}{%
\section{The 5 \ensuremath{\times } 5 test-asset grid}\label{the-5-5-test-asset-grid}}

At each annual formation date, every eligible US common stock
(NYSE, AMEX and NASDAQ) is placed into one of five \emph{size} groups by
market equity and, independently, one of five \emph{book-to-market}
groups. The $5\times5=25$ intersections define 25 portfolios whose monthly
returns are \emph{value-weighted}: each constituent enters in proportion
to its market equity, so a handful of large firms can dominate a cell. The
breakpoints are recomputed each year from NYSE stocks only, and a stock
stays in its assigned cell until the next formation date. A label such as
\texttt{SMALL\ HiBM} therefore denotes a \emph{characteristic cell} ---
small firms with a high book-to-market ratio --- not a company, a managed
fund, or a trading recommendation. The corner cells name the extremes:
\texttt{SMALL\ LoBM} is small-growth, \texttt{SMALL\ HiBM} is small-value,
\texttt{BIG\ LoBM} is large-growth and \texttt{BIG\ HiBM} is large-value;
the interior cells are written \texttt{ME}$j$~\texttt{BM}$k$ for size
quintile $j$ and book-to-market quintile $k$.

Two properties of this grid drive design choices later. First, the cells
are \emph{not independent entities}: neighbours share construction rules
and, through the common market, share monthly shocks --- the source of the
cross-sectional dependence handled by the joint covariance of
Chapter~\ref{preliminaries} (\S\ref{prelim-hac}). Second, the grid is
built from the very sorts (size and book-to-market) that define the SMB
and HML factors, so regressing these portfolios on FF3 tests the model
against its own raw material (\S\ref{prelim-ff}).

\begin{longtable}[]{@{}p{0.20\textwidth}p{0.34\textwidth}p{0.36\textwidth}@{}}
\toprule\noalign{}
Term & Meaning & Common mistake \\
\midrule\noalign{}
\endhead
\bottomrule\noalign{}
\endlastfoot
Size & Market equity grouping & Confusing ``small'' with low price
rather than low market capitalisation \\
Book-to-market & Book equity divided by market equity & Reversing value
and growth: high B/M is value; low B/M is growth \\
Value-weighted return & Constituent returns weighted by market value &
Accidentally parsing the equal-weighted block \\
Portfolio-month & One portfolio observed in one calendar month & Joining
files many-to-many and duplicating this grain \\
\end{longtable}

\hypertarget{returns-and-factors}{%
\section{Returns and factors}\label{returns-and-factors}}

The dependent variable of every regression is the portfolio's monthly
\emph{excess} return,
\begin{equation}
y_{it} = r_{it} - r_{ft},
\end{equation}
the total return $r_{it}$ of portfolio $i$ in month $t$ minus the
one-month risk-free rate $r_{ft}$. The regressors are the traded factor
returns collected in $f_t$:
\begin{itemize}
\tightlist
\item
  \textbf{Mkt-RF:} the aggregate market return in excess of the risk-free
  rate (itself an excess return).
\item
  \textbf{SMB} (small minus big): the size factor, long small firms and
  short large ones.
\item
  \textbf{HML} (high minus low): the value factor, long high
  book-to-market and short low.
\item
  \textbf{RMW and CMA:} profitability (robust minus weak) and investment
  (conservative minus aggressive) factors, added by the five-factor model.
\end{itemize}

Because each regressor is itself an excess or self-financing long--short
return, the dependent variable and the regressors sit on the same
economic footing, and the intercept $\alpha_i$ inherits the units of a
monthly excess return (\S\ref{prelim-ap}).

\textbf{Units and annualisation.} The canonical analysis works in
\emph{decimal monthly} returns, so a one-percent monthly return is
$0.01$, not $1$. A monthly decimal alpha $a$ is reported in annualised
basis points as $a\times 12\times 10{,}000$; for example $a=0.001$ becomes
$0.001\times 12\times 10{,}000 = 120$ bps/year. This is a \emph{linear}
reporting convention, not geometric compounding, chosen so that the
reported figure is a transparent scaling of the estimated coefficient.
Mixing percent and decimal units is the most common unit error and would
inflate every alpha and standard error by a factor of one hundred; the
loader therefore records and checks the units explicitly.

\hypertarget{finance-foundations-for-a-mathematics-graduate}{%
\section{Finance foundations for a mathematics
graduate}\label{finance-foundations-for-a-mathematics-graduate}}

\hypertarget{simple-and-portfolio-returns}{%
\subsection{Simple and portfolio
returns}\label{simple-and-portfolio-returns}}

If an asset begins the month at price P\textsubscript{t\ensuremath{-}1}, pays
dividend D\textsubscript{t}, and ends at P\textsubscript{t}, its simple
total return is

\begin{equation}
R_{t} = (P_{t} + D_{t} - P_{t-1}) / P_{t-1}.
\end{equation}

For J constituents with beginning-of-period portfolio weights
w\textsubscript{jt\ensuremath{-}1}, \ensuremath{\sum }\textsubscript{j}w\textsubscript{jt\ensuremath{-}1}=1, the
one-period portfolio return is

\begin{equation}
R_{\mathrm{pt}} = \sum _{j=1}^{J} w_{jt-1}R_{\mathrm{jt}}.
\end{equation}

In a value-weighted portfolio, w\textsubscript{jt\ensuremath{-}1} is proportional to
the firm\textquotesingle s market equity at the weighting date. Larger
companies therefore receive larger weights.

\hypertarget{why-subtract-the-risk-free-rate}{%
\subsection{Why subtract the risk-free
rate?}\label{why-subtract-the-risk-free-rate}}

A risky investment must first compensate the investor for deferring
consumption. The risk-free return r\textsubscript{ft} is the baseline
return available without bearing the portfolio\textquotesingle s risky
exposure. The excess return
y\textsubscript{it}=r\textsubscript{it}\ensuremath{-}r\textsubscript{ft} is the part
to be explained by risky factors and alpha.

\hypertarget{market-equity-and-book-to-market}{%
\subsection{Market equity and
book-to-market}\label{market-equity-and-book-to-market}}

\begin{equation}
\text{Market equity} = \text{share price} \times \text{shares outstanding},
\end{equation}
\begin{equation}
\text{Book-to-market} = \frac{\text{accounting book equity}}{\text{market equity}}.
\end{equation}

Market equity (``size'') measures how large the firm is in the market;
book-to-market compares the accountant's valuation of the firm's equity to
the market's. High book-to-market companies are conventionally called
\emph{value} stocks --- the market prices them low relative to book --- and
low book-to-market companies are called \emph{growth} stocks. These labels
describe \emph{characteristics}, not guaranteed future returns, and the
value/growth direction is the single most common labelling error: high
B/M is value, low B/M is growth.

\hypertarget{how-the-ff3-factor-portfolios-are-constructed}{%
\subsection{How the FF3 factor portfolios are
constructed}\label{how-the-ff3-factor-portfolios-are-constructed}}

The 25 portfolios in this study are $5\times5$ \textbf{test assets}. They
are related to, but not identical with, the six $2\times3$ portfolios used
in the standard construction of the SMB and HML \emph{factors}. Writing
$S$/$B$ for the two size groups (small, big) and $L$/$N$/$H$ for the three
book-to-market groups (low, neutral, high), the six building-block
portfolios are the intersections, and the textbook factor formulas average
them:
\begin{equation}
\mathrm{SMB} = \tfrac13(S/L + S/N + S/H) - \tfrac13(B/L + B/N + B/H),
\end{equation}
\begin{equation}
\mathrm{HML} = \tfrac12(S/H + B/H) - \tfrac12(S/L + B/L).
\end{equation}
SMB is thus the average small-cap return minus the average large-cap
return holding book-to-market spread fixed, and HML is the average
high-B/M return minus the average low-B/M return holding size fixed. Both
are \emph{self-financing long--short} returns --- the long and short legs
carry offsetting capital weights --- so they are returns on zero-cost
portfolios, and Mkt-RF is likewise an excess return. The five-factor model
appends two more long--short returns built the same way,
\begin{equation}
\mathrm{RMW} = \text{(robust-profitability)} - \text{(weak-profitability)},
\end{equation}
\begin{equation}
\mathrm{CMA} = \text{(conservative-investment)} - \text{(aggressive-investment)}.
\end{equation}

The distinction between the $2\times3$ factor-building portfolios and the
$5\times5$ test assets matters: the study regresses the \emph{test assets}
on the \emph{factors}, and because both come from size and book-to-market
sorts the resulting alpha is a residual of the model against its own
inputs.

\textbf{Economic meaning of the regression.} The regression
$y_{it}=\alpha_i+\beta_i^{\top}f_t+\varepsilon_{it}$ asks how much of a
portfolio's excess return moves one-for-one with the broad market, size
and value returns; the loadings $\beta_i$ are those sensitivities and the
intercept $\alpha_i$ is the average return left after the fitted linear
co-movement is removed. A positive $\beta_{i,\mathrm{HML}}$, for instance,
says the cell behaves like a value portfolio, which the sanity checks of
Chapter~13 expect to rise across the book-to-market columns.

\hypertarget{fail-closed-panel-validation}{%
\section{Fail-closed panel
validation}\label{fail-closed-panel-validation}}

The analytical grain is exactly one \emph{portfolio-month}: one row per
portfolio per calendar month, $25\times1{,}193 = 29{,}825$ rows in the
canonical panel. A loader is \emph{fail-closed} when it raises an error
rather than silently proceeding on any violation of that grain, so that a
corrupt panel can never reach a regression. Six checks are enforced, each
guarding against a specific way inference can be quietly invalidated:
\begin{itemize}
\tightlist
\item
  \textbf{Required columns present} --- a missing field (say the risk-free
  rate) would otherwise be filled with defaults or \texttt{NaN} downstream.
\item
  \textbf{No null required values} --- a single \texttt{NaN} propagates
  through a regression and can silently drop observations.
\item
  \textbf{No duplicate (portfolio, date) keys} --- duplicates double-count
  months and corrupt standard errors, likelihoods and ranks.
\item
  \textbf{Balanced panel} --- every month must contain all 25 portfolios,
  because the joint HAC covariance and the common-date bootstrap require a
  rectangular time-by-portfolio array.
\item
  \textbf{No missing calendar month} --- a gap in the monthly sequence
  would break the lag structure of the HAC estimator (\S\ref{prelim-hac}).
\item
  \textbf{Within-month factor agreement} --- the factor values are common
  to all portfolios in a month, so any within-month disagreement signals a
  bad join.
\end{itemize}

\begin{verbatim}
REQUIRED = {"date", "portfolio", "return", "mkt_rf", "smb", "hml", "rf"}

def validate_panel(df, expected_portfolios=25):
    missing = REQUIRED - set(df.columns)
    if missing:
        raise ValueError(f"Missing columns: {sorted(missing)}")
    if df[list(REQUIRED)].isna().any().any():
        raise ValueError("Null required values")
    if df.duplicated(["portfolio", "date"]).any():
        raise ValueError("Duplicate portfolio-month keys")

    counts = df.groupby("date")["portfolio"].nunique()
    if not counts.eq(expected_portfolios).all():
        raise ValueError("Unbalanced panel")

    dates = pd.PeriodIndex(sorted(df["date"].unique()), freq="M")
    expected = pd.period_range(dates.min(), dates.max(), freq="M")
    if not dates.equals(expected):
        raise ValueError("Missing calendar month")

    factor_spread = df.groupby("date")[["mkt_rf","smb","hml","rf"]].nunique()
    if factor_spread.to_numpy().max() != 1:
        raise ValueError("Factor values disagree within month")
\end{verbatim}

\textbf{A many-to-many join is not more data.} If the same
portfolio-month appears more than once after a merge, standard errors,
likelihoods and rankings can all become invalid. Always validate the
analytical key immediately after joining.

\hypertarget{provenance}{%
\section{Provenance and the replication/reconciliation distinction}\label{provenance}}

Reproducibility requires knowing exactly which bytes produced a result.
For each input the pipeline records the source URL, the acquisition date
if known, the sample end date, the file size, a SHA-256 content hash, the
weighting block used, the units, and the parser version. A SHA-256 digest
is a fixed-length fingerprint of the file's bytes: identical files share a
digest, and any change almost certainly changes it, so a recorded hash
lets a later reader confirm they hold the same file. Note that a sample
\emph{ending} in November 2025 does not prove the file was
\emph{downloaded} in November 2025 --- the coverage date and the
acquisition vintage are different facts.

\begin{verbatim}
import hashlib

def sha256(path):
    h = hashlib.sha256()
    with open(path, "rb") as f:
        for chunk in iter(lambda: f.read(1 << 20), b""):
            h.update(chunk)
    return h.hexdigest()
\end{verbatim}

\textbf{Replication is not reconciliation.} Two different guarantees are
often confused. \emph{Numerical replication} asks: given \emph{these}
bytes, does an independent implementation reproduce the reported
coefficients? Here it does, to a tolerance of $10^{-12}$ against
\texttt{statsmodels} (Chapter~15). \emph{Source reconciliation} asks: are
\emph{these} bytes the same as the official public source? That is a
separate question, and it does not pass.

\textbf{Paper-specific limitation.} The fixed local inputs and the current
(May 2026) official snapshot from the Kenneth R.\ French Data Library
differ over their overlapping period --- by up to $1.4167$ percentage
points for portfolio returns and $0.34$ percentage points for the FF3/RF
fields. Official historical revisions are a plausible but unproven cause.
Because the local acquisition vintage is unverified, exact source equality
cannot be asserted and is classified \textsc{not\_comparable} rather than
\textsc{pass} --- a visible, non-blocking finding, not a silent one --- even
though every within-pipeline integrity check above passes. The local
inputs remain immutable and are identified by hash, so the analysis is
fully reproducible even while its reconciliation to the live source
remains open (Chapter~15).

